\documentclass[CJK,13pt]{beamer}
\usepackage{CJKutf8}
\usepackage{beamerthemesplit}
\usetheme{Malmoe}
\useoutertheme[footline=authortitle]{miniframes}
\usepackage{amsmath}
\usepackage{amssymb}
\usepackage{graphicx}
\usepackage{eufrak}
\usepackage{color}
\usepackage{slashed}
\usepackage{simplewick}
\usepackage{tikz}
\usepackage{tcolorbox}
\usepackage{ulem}
\graphicspath{{../figures/}}
%%figures
\def\lfig#1#2{\includegraphics[width=#1 in]{#2}}
\def\tfig#1#2{\includegraphics[height=#1 in]{#2}}
\def\addfig#1#2{\begin{center}\includegraphics[width=#1 in]{#2}\end{center}}
\def\question{\includegraphics[width=0.3in]{why.jpg}\,}
\def\answer{\includegraphics[width=0.3in]{baozou_haha.png}\,}
\def\wulian{\includegraphics[width=0.18in]{emoji_wulian.jpg}}
\def\bigwulian{\includegraphics[width=0.35in]{emoji_wulian.jpg}}
\def\bye{\includegraphics[width=0.18in]{emoji_bye.jpg}}
\def\bigbye{\includegraphics[width=0.35in]{emoji_bye.jpg}}
\def\huaixiao{\includegraphics[width=0.18in]{emoji_huaixiao.jpg}}
\def\bighuaixiao{\includegraphics[width=0.35in]{emoji_huaixiao.jpg}}
\def\jianxiao{\includegraphics[width=0.18in]{emoji_jianxiao.jpg}}
\def\bigjianxiao{\includegraphics[width=0.35in]{emoji_jianxiao.jpg}}
\def\haoqi{\includegraphics[width=0.18in]{emoji_haoqi.jpg}}
%% colors
\def\blacktext#1{{\color{black}#1}}
\def\bluetext#1{{\color{blue}#1}}
\def\redtext#1{{\color{red}#1}}
\def\darkbluetext#1{{\color[rgb]{0,0.2,0.6}#1}}
\def\skybluetext#1{{\color[rgb]{0.2,0.7,1.}#1}}
\def\cyantext#1{{\color[rgb]{0.,0.5,0.5}#1}}
\def\greentext#1{{\color[rgb]{0,0.7,0.1}#1}}
\def\darkgray{\color[rgb]{0.2,0.2,0.2}}
\def\lightgray{\color[rgb]{0.6,0.6,0.6}}
\def\gray{\color[rgb]{0.4,0.4,0.4}}
\def\blue{\color{blue}}
\def\red{\color{red}}
\def\orange{\color[rgb]{1.,0.8,0.}}
\def\green{\color{green}}
\def\darkgreen{\color[rgb]{0,0.4,0.1}}
\def\darkblue{\color[rgb]{0,0.2,0.6}}
\def\skyblue{\color[rgb]{0.2,0.7,1.}}
%%control
\def\heaviside{\,\mathrm{h}\,}
\def\bral{\left(\begin{array}{l}}
\def\brar{\end{array}\right)}
\def\brall{\left(\begin{array}{ll}}
\def\brarr{\end{array}\right)}
\def\bralll{\left(\begin{array}{lll}}
\def\brarrr{\end{array}\right)}
\def\branchl{\left\{\begin{array}{l}}
\def\branchr{\end{array}\right.}
\def\branchll{\left\{\begin{array}{ll}}
\def\branchrr{\end{array}\right.}
\def\branchlll{\left\{\begin{array}{lll}}
\def\branchrrr{\end{array}\right.}
\def\sfgamma#1{\,\Gamma\left( #1 \right)\,}
\def\be{\begin{equation}}
\def\ee{\nonumber\end{equation}}
\def\bea{\begin{eqnarray}}
\def\eea{\nonumber\end{eqnarray}}
\def\bch{\begin{CJK}{UTF8}{gbsn}}
\def\ech{\end{CJK}}
\def\bitem{\begin{itemize}}
\def\eitem{\end{itemize}}
\def\bcenter{\begin{center}}
\def\ecenter{\end{center}}
\def\bex{\begin{minipage}{0.2\textwidth}\includegraphics[width=0.6in]{jugelizi.png}\end{minipage}\begin{minipage}{0.76\textwidth}}
\def\eex{\end{minipage}}
\def\chtitle#1{\frametitle{\bch#1\ech}}
\def\bmat#1{\left(\begin{array}{#1}}
\def\emat{\end{array}\right)}
\def\bcase#1{\left\{\begin{array}{#1}}
\def\ecase{\end{array}\right.}
\def\bmini#1{\begin{minipage}{#1\textwidth}}
\def\emini{\end{minipage}}
\def\tbox#1{\begin{tcolorbox}#1\end{tcolorbox}}
\def\pfrac#1#2#3{\left(\frac{\partial #1}{\partial #2}\right)_{#3}}
\def\res#1#2{\,\mathrm{res}\,#1\left(#2\right)\,}
\def\newt#1#2{\left(\begin{array}{c}#1\\ #2\end{array}\right)}
\def\reof#1{\,\mathrm{Re}{\left(#1\right)\,}}
\def\imof#1{\,\mathrm{Im}{\left(#1\right)\,}}
\def\Arg#1{\,\mathrm{Arg}\,#1\,}
%%symbols
\def\intfull{\int_{-\infty}^\infty}
\def\inthalf{\int_0^\infty}
\def\bropt{\,(\ \ \ )}
\def\sech{\mathrm{sech}\,}
\def\csch{\mathrm{csch}\,}
\def\asinh{\mathrm{asinh}\,}
\def\acosh{\mathrm{acosh}\,}
\def\sone{$\star$}
\def\stwo{$\star\star$}
\def\sthree{$\star\star\star$}
\def\sfour{$\star\star\star\star$}
\def\sfive{$\star\star\star\star\star$}
\def\rint{{\int_\leftrightarrow}}
\def\roint{{\oint_\leftrightarrow}}
\def\stdHf{{\textit{\r H}_f}}
\def\deltaH{{\Delta \textit{\r H}}}
\def\ii{{\dot{\imath}}}
\def\skipline{{\vskip0.1in}}
\def\skiplines{{\vskip0.2in}}
\def\lagr{{\mathcal{L}}}
\def\hamil{{\mathcal{H}}}
\def\vecT{{\mathbf{T}}}
\def\vecN{{\mathbf{N}}}
\def\vecB{{\mathbf{B}}}
\def\vecv{{\mathbf{v}}}
\def\vecr{{\mathbf{r}}}
\def\vecf{{\mathbf{f}}}
\def\vecg{{\mathbf{g}}}
\def\vecupsilon{{\mathbf{\upsilon}}}
\def\vecu{{\mathbf{u}}}
\def\vecj{{\mathbf{j}}}
\def\vecx{{\mathbf{x}}}
\def\vecy{{\mathbf{y}}}
\def\vecz{{\mathbf{z}}}
\def\veck{{\mathbf{k}}}
\def\vecp{{\mathbf{p}}}
\def\vecn{{\mathbf{n}}}
\def\vecA{{\mathbf{A}}}
\def\vecP{{\mathbf{P}}}
\def\vecsigma{{\mathbf{\sigma}}}
\def\hatJn{{\hat{J_\vecn}}}
\def\hatJx{{\hat{J_x}}}
\def\hatJy{{\hat{J_y}}}
\def\hatJz{{\hat{J_z}}}
\def\hatj#1{\hat{J_{#1}}}
\def\hatphi{{\hat{\phi}}}
\def\hatq{{\hat{q}}}
\def\hatpi{{\hat{\pi}}}
\def\vel{\upsilon}
\def\Dint{{\mathcal{D}}}
\def\adag{{\hat{a}^\dagger}}
\def\bdag{{\hat{b}^\dagger}}
\def\cdag{{\hat{c}^\dagger}}
\def\ddag{{\hat{d}^\dagger}}
\def\hata{{\hat{a}}}
\def\hatb{{\hat{b}}}
\def\hatc{{\hat{c}}}
\def\hatd{{\hat{d}}}
\def\hatD{{\,\hat{D}}}
\def\hatN{{\hat{N}}}
\def\hatH{{\hat{H}}}
\def\hatp{{\hat{p}}}
\def\Fup{{F^{\mu\nu}}}
\def\Fdown{{F_{\mu\nu}}}
\def\newl{\nonumber \\}
\def\vece{\mathrm{e}}
\def\calM{{\mathcal{M}}}
\def\calT{{\mathcal{T}}}
\def\calR{{\mathcal{R}}}
\def\barpsi{\bar{\psi}}
\def\baru{\bar{u}}
\def\barv{\bar{\upsilon}}
\def\qeq{\stackrel{?}{=}}
\def\ftf{\stackrel{\mathcal{FT}}{\Longrightarrow}}
\def\ftb{\stackrel{\mathcal{FT}}{\Longleftarrow}}
\def\ftfb{\stackrel{\mathcal{FT}}{\Longleftrightarrow}}
\def\ltf{\stackrel{\mathcal{LT}}{\Longrightarrow}}
\def\ltb{\stackrel{\mathcal{LT}}{\Longleftarrow}}
\def\ltfb{\stackrel{\mathcal{LT}}{\Longleftrightarrow}}
\def\torder#1{\mathcal{T}\left(#1\right)}
\def\rorder#1{\mathcal{R}\left(#1\right)}
\def\contr#1#2{\contraction{}{#1}{}{#2}#1#2}
\def\trof#1{\mathrm{Tr}\left(#1\right)}
\def\trace{\mathrm{Tr}}
\def\comm#1{\ \ \ \left(\mathrm{used}\ #1\right)}
\def\tcomm#1{\ \ \ (\text{#1})}
\def\slp{\slashed{p}}
\def\slk{\slashed{k}}
\def\calp{{\mathfrak{p}}}
\def\veccalp{\mathbf{\mathfrak{p}}}
\def\Tthree{T_{\tiny \textcircled{3}}}
\def\pthree{p_{\tiny \textcircled{3}}}
\def\dbar{{\,\mathchar'26\mkern-12mu d}}
\def\erf{\mathrm{erf}}
\def\const{\mathrm{const.}}
\def\pheat{\pfrac p{\ln T}V}
\def\vheat{\pfrac V{\ln T}p}
%%units
\def\fdeg{{^\circ \mathrm{F}}}
\def\cdeg{^\circ \mathrm{C}}
\def\atm{\,\mathrm{atm}}
\def\angstrom{\,\text{\AA}}
\def\SIL{\,\mathrm{L}}
\def\SIkm{\,\mathrm{km}}
\def\SIyr{\,\mathrm{yr}}
\def\SIGyr{\,\mathrm{Gyr}}
\def\SIV{\,\mathrm{V}}
\def\SImV{\,\mathrm{mV}}
\def\SIeV{\,\mathrm{eV}}
\def\SIkeV{\,\mathrm{keV}}
\def\SIMeV{\,\mathrm{MeV}}
\def\SIGeV{\,\mathrm{GeV}}
\def\SIcal{\,\mathrm{cal}}
\def\SIkcal{\,\mathrm{kcal}}
\def\SImol{\,\mathrm{mol}}
\def\SIN{\,\mathrm{N}}
\def\SIHz{\,\mathrm{Hz}}
\def\SIm{\,\mathrm{m}}
\def\SIcm{\,\mathrm{cm}}
\def\SIfm{\,\mathrm{fm}}
\def\SImm{\,\mathrm{mm}}
\def\SInm{\,\mathrm{nm}}
\def\SImum{\,\mathrm{\mu m}}
\def\SIJ{\,\mathrm{J}}
\def\SIW{\,\mathrm{W}}
\def\SIkJ{\,\mathrm{kJ}}
\def\SIs{\,\mathrm{s}}
\def\SIkg{\,\mathrm{kg}}
\def\SIg{\,\mathrm{g}}
\def\SIK{\,\mathrm{K}}
\def\SImmHg{\,\mathrm{mmHg}}
\def\SIPa{\,\mathrm{Pa}}
%page
\def\secpage#1#2{\begin{frame}\bcenter{\bf \Huge #1} \ecenter \skipline \tbox{\bcenter #2 \ecenter}\end{frame}}
\def\append#1#2{\secpage{附录#1}{#2}}
\def\thinka#1{\begin{frame}\frametitle{思考题}\bcenter\lfig{0.4}{think0.jpg}  \skipline #1 \ecenter\end{frame}}
\def\thinkb#1{\begin{frame}\frametitle{思考题}\bcenter\lfig{0.5}{think1.jpg}  \skipline #1 \ecenter\end{frame}}
\def\thinkc#1{\begin{frame}\frametitle{思考题}\bcenter\lfig{0.5}{think2.jpg}  \skipline #1 \ecenter\end{frame}}
\def\thinkd#1{\begin{frame}\frametitle{思考题}\bcenter\lfig{0.5}{think3.jpg}  \skipline #1 \ecenter\end{frame}}
\def\thinke#1{\begin{frame}\frametitle{思考题}\bcenter\lfig{0.5}{think4.jpg}  \skipline #1 \ecenter\end{frame}}
\def\thinkf#1{\begin{frame}\frametitle{思考题}\bcenter\lfig{0.5}{think5.jpg}  \skipline #1 \ecenter\end{frame}}

\def\courseurl{http://zhiqihuang.top/gr}

\def\tpage#1#2{
\title{GR \S{#1}  #2}
  \author{Zhiqi Huang}

\begin{frame}
\begin{center}
{\bf \Huge G}eneral {\bf \Huge R}elativity

{\vskip 0.1in}



{\Large \S #1 #2}

{\vskip 0.2in}

{Lecturer: 黄志琦}

\vskip 0.2in

\courseurl

\end{center}
\end{frame}
}


  \date{}
  \begin{document}
  \bch
\tpage{5}{Connection}



\begin{frame}
  回忆一下上一讲的结果: 设三维欧氏空间里的曲面 $\vecx(u^1,u^2)$ 有第一基本形式
  $$ds^2 = g_{ij}du^idu^j,$$
  和第二基本形式
  $$\vecn\cdot d^2\vecx = \beta_{ij}du^i du^j,$$  
  这里的 $\vecn$ 是法向量， $g_{ij} = \vecx_{,i}\cdot\vecx_{,j}$, $\beta_{ij} = \vecn \cdot \vecx_{,i,j} = -\vecn_{,i} \cdot \vecx_{,j} = -\vecx_{,i} \cdot \vecn_{,j}$。

  
  本讲的任务是用第一基本形式（度规 $g_{ij}$）计算高斯曲率 $K = \frac{\det(\beta)}{\det(g)}$。

  \skiplines
  
  在此之前，我们又要引入一些新的符号：请准备好呼吸机……
\end{frame}

\begin{frame}
  \frametitle{度规和指标的升降}
 第一基本形式里的度规矩阵
  \be
  \begin{pmatrix}
    g_{11} & g_{12} \\
    g_{21} & g_{22} 
  \end{pmatrix}
  \ee
  由高维欧式空间的距离平方导出，所以一定是正定矩阵。但是在更一般的讨论中，我们将放宽这一假设，只要求它是满秩矩阵。

  不管怎么说，度规矩阵的逆矩阵存在，我们把它写成
  \be
  \begin{pmatrix}
    g^{11} & g^{12} \\
    g^{21} & g^{22} 
  \end{pmatrix}
  \equiv   \begin{pmatrix}
    g_{11} & g_{12} \\
    g_{21} & g_{22} 
  \end{pmatrix}^{-1}
  \ee
  度规矩阵和它的逆的乘积当然是单位矩阵，这一事实用爱因斯坦的求和约定写出来就是
  \be
  g_{ij}g^{jk} = \delta_i^{\ k} \equiv \left\{
  \begin{matrix}
    1, & \text{if } i = k \\
    0, & \text{else}
  \end{matrix}
  \right.
  \ee
\end{frame}


\begin{frame}
  \frametitle{指标的升降}
  我们搞那么多事情，一会儿把指标写上面(例如 $du^i$)，一会儿把指标写下面（例如 $g_{ij}$），是因为上下指标含义不同，且可以互相变换。

  \skipline
  
 具体来说，我们将遵循这样的原则对带指标的量进行“指标升降”。
  \bitem
\item{如果要把一个上指标 $\alpha$ 降为下指标，我们把上指标里的 $\alpha$ 换成任何一个新的没有用到过的指标，例如 $\beta$，然后乘以 $g_{\beta\alpha}$。}
\item{如果要把一个下指标 $\alpha$ 升为上指标，我们把下指标里的 $\alpha$ 换成任何一个新的没有用到过的指标，例如 $\beta$，然后乘以 $g^{\beta\alpha}$。}
  \eitem
  当然，上面规则中的$\beta$ 这个指标是什么不重要，因为按爱因斯坦求和规则我们会对 $\beta$ 进行遍历求和。  
\end{frame}


\begin{frame}
  例如，我们可以把 $du^i$ 里的上指标 $i$ 降下去，具体就是：
  $$ du_i \equiv g_{ij} du^j.$$
  粗略地看，这无非就是个定义而已。

  \skipline

  问题是，这样的定义是自洽的吗？如果我们把 $du_i$ 里的下指标再升上去，还会是和原来一样的 $du^i$ 吗？
\end{frame}



\begin{frame}
  答案是肯定的：

  $$ g^{ij}du_j = g^{ij} g_{jk} du^k = \delta^i_{\ k}du^k = du^i.$$

  很容易把这个简短的证明推广到任意多个指标的任意升降。
\end{frame}


\thinkf{$g^{ij}$ 比较特殊，它是通过 $g_{ij}$矩阵的逆矩阵定义的。请检验,
  \bitem
\item{如果把 $g_{ij}$ 的指标按照上述升降指标的规则升上去，会得到（按逆矩阵定义的） $g^{ij}$ 吗？ }
\item{仅上升一个指标的度规 $g_i^{\ j}$ 或者 $g^i_{\ j}$ 是什么？}  
  \eitem}


\begin{frame}
  \frametitle{指标升降操和偏导操作不可交换}

  对坐标的偏导会产生额外的指标。要格外小心的是，指标的升降操作和偏导产生新指标的操作一般是不可交换的。例如： 我们可以把 $g_{ij}$ 的指标 $i$ 升上去得到 $g^i_{\ j}$，然后求偏导得到 $g^i_{\ j, k}$。如果我们先求偏导，得到 $g_{ij, k}$ ，然后把指标 $i$ 升上去，会得到一个数值上不一样的结果。

  \skiplines

  在我们之后学习了张量和协变导数之后，这个问题会得到解决（协变导数和指标升降的操作可以交换）。目前我们对带着偏导的量的指标还是要十分小心，不要随便进行指标升降。

\end{frame}


\begin{frame}
  \frametitle{小心地应用指标升降的例子}
  \addfig{0.8}{eg.jpg}
  
  虽然 $\vecx_{,i}$ 带着偏导指标，我们还是可以强行定义 $\vecx^{,i}=g^{ij}\vecx_{,j}$。由于 $\vecx^{,i}$ 并没有其他什么定义方式，这种操作还不至于导致混乱。

  \skiplines
  
  容易严格按照定义验证
  \be
  \vecx_{,i}\cdot \vecx^{,j} = g_i^{\ j}=\left\{
  \begin{array}{ll}    1, & \text{ if } i=j \\
    0, & \text{ else}
  \end{array}
  \right.
  \ee
  
  $$\vecx^{,i}\vecx^{,j} = g^{ij}$$

\end{frame}




\begin{frame}
  好了，摘下呼吸机。我们要开始操作了。
\end{frame}


\begin{frame}  

  \frametitle{秘密武器}
  我们先来介绍一个秘密武器：{\blue  三维欧氏空间任意矢量 $\vecv$ 可以分解成
  $$ \vecv = \left(\vecx_{,i}\cdot\vecv\right)\vecx^{,i}  + (\vecn\cdot\vecv)\vecn.$$}
  这个等式是这样证明的，在两边点乘 $\vecx_{,j}$ 或 $\vecn$ 都得到相同的结果（这里需要用到 $\vecx_{,i}\cdot\vecx^{,j} = g_i^{\ j}$）。由于 $\vecx_{,1}$, $\vecx_{,2}$, $\vecn$ 是完备的，故得证。

  \addfig{1}{caozuo.jpg}

\end{frame}

\begin{frame}
  \frametitle{联络(connection)}
  我们定义一个非常有用的辅助记号： {\blue 联络
    $$\Gamma_{kij} \equiv \vecx_{,k}\cdot\vecx_{,i,j}.$$}
  显然这样定义的联络具有 $\Gamma_{kij} = \Gamma_{kji}$ 的对称性。

  \skiplines
  
    我们还经常需要把联络 $\Gamma_{kij}$ 的第一个指标升上去，变成
  {\blue
  $$\Gamma^k_{\ ij} \equiv g^{kl}\Gamma_{lij} = \vecx^{,k}\cdot\vecx_{,i,j}$$}

\end{frame}


\begin{frame}
  \frametitle{联络可以用度规导出}
  度规的偏导数可以用联络写出来：
  $$g_{ij,k} = \left(\vecx_{,i}\cdot\vecx_{,j}\right)_{,k} = \vecx_{,i,k}\cdot\vecx_{,j} + \vecx_{,i}\cdot\vecx_{,j, k}  = \Gamma_{jik}+\Gamma_{ijk}.$$
  同理
  $$ g_{ki,j} = \Gamma_{ijk}+\Gamma_{kij}$$
  $$ g_{kj,i} = \Gamma_{jik}+\Gamma_{kij}$$
  把后两式相加并减去第一式，即可得到：
  \tbox{$$\Gamma_{kij} = \frac{1}{2}\left(g_{ki,j}+g_{kj,i}-g_{ij,k}\right).$$}
  也就是说，联络可以由第一基本形式得到。
\end{frame}


\begin{frame}
  \frametitle{升级武器}
  我们再把秘密武器进行一下升级：任意两个三维欧式空间的矢量 $\vecy$, $\vecz$。可以把它们各自按秘密武器展开后，再求两者的内积
  \bea
  \vecy\cdot\vecz &=& \left[\left(\vecx_{,i}\cdot\vecy\right)\vecx^{,i}+(\vecn\cdot\vecy)\vecn\right]\cdot\left[\left(\vecx_{,j}\cdot\vecz\right)\vecx^{,j}+(\vecn\cdot\vecz)\vecn\right] \newl
  &=& \left(\vecx_{,i}\cdot\vecy\right)\left(\vecx_{,j}\cdot\vecz\right)\vecx^{,i}\cdot\vecx^{,j}  + (\vecn\cdot\vecy)(\vecn\cdot\vecz) \newl
  &=& \left(\vecx_{,i}\cdot\vecy\right)\left(\vecx_{,j}\cdot\vecz\right)g^{ij} + (\vecn\cdot\vecy)(\vecn\cdot\vecz) \newl
  &=& \left(\vecx_{,i}\cdot\vecy\right)\left(\vecx^{,i}\cdot\vecz\right) + (\vecn\cdot\vecy)(\vecn\cdot\vecz)    
  \eea
  我们得到了一个计算 $(\vecn\cdot y)(\vecn\cdot z)$ 的有效方法：
  {\blue
     $$ (\vecn\cdot\vecy)(\vecn\cdot\vecz)   =  \vecy\cdot\vecz - \left(\vecx_{,i}\cdot\vecy\right)\left(\vecx^{,i}\cdot\vecz\right).$$
  }

\end{frame}



\begin{frame}
  \frametitle{计算高斯曲率 $K=\frac{\det(\beta)}{\det(g)}$}
  利用升级的秘密武器，立刻有
 {\small
  \bea
  && \det(\beta) \newl
  &=& \beta_{11}\beta_{22} - (\beta_{12})^2 \newl
  &=& \left(\vecn\cdot\vecx_{,1,1}\right)\left(\vecn\cdot\vecx_{,2,2}\right) - \left(\vecn\cdot\vecx_{,1,2}\right)^2 \newl
  &=& \vecx_{,1,1}\cdot\vecx_{,2,2} - \left(\vecx_{,i}\cdot\vecx_{,1,1}\right)\left(\vecx^{,i}\cdot\vecx_{,2,2}\right) - \left(\vecx_{,1,2}\right)^2 +  \left(\vecx_{,i}\cdot\vecx_{,1,2}\right)\left(\vecx^{,i}\cdot\vecx_{,1,2}\right) \newl
  &=& \vecx_{,1,1}\cdot\vecx_{,2,2}  - \Gamma_{i11}\Gamma^i_{\ 22} -  \left(\vecx_{,1,2}\right)^2 + \Gamma_{i12}\Gamma^i_{\ 12} \newl 
  &=& (\vecx_{,1,1}\cdot\vecx_{,2})_{,2}-  \left(\vecx_{,1,2}\cdot\vecx_{,2}\right)_{,1} + \Gamma_{i12}\Gamma^i_{\ 12} - \Gamma_{i11}\Gamma^i_{\ 22} \newl
  &=& \left(\Gamma_{211}\right)_{,2} -  \left(\Gamma_{212}\right)_{,1} + \Gamma_{i12}\Gamma^i_{\ 12} - \Gamma_{i11}\Gamma^i_{\ 22} \newl    
 \eea
 }
 联络可以用度规和度规的偏导数写出来，所以任务完成。
\end{frame}


\begin{frame}
  \frametitle{sympy又来了}
  从度规到联络和从联络到高斯曲率的计算都比较繁琐。我们可以写一个 sympy 脚本代劳，请参考：

  \url{http://zhiqihuang.top/gr/codes/connection.py}
\end{frame}


\ech
\end{document}
